problem:  
Let \((M,\Delta,g)\) be a **compact equiregular sub-Riemannian manifold** with bracket-generating distribution \(\Delta \subset TM\) and sub-Riemannian metric \(g\).  
Denote by \(\Delta^\perp \subset T^*M\) the annihilator of \(\Delta\). For each \(\lambda \in \Delta^\perp\), define the **abnormal Hamiltonian system** associated with smooth vector fields \(X_1, \dots, X_k\) spanning \(\Delta\):  
\[
\dot{\lambda}(t) = \sum_{i=1}^k u_i(t)\,\vec{h}_i(\lambda(t)), \quad \text{where } h_i(\lambda) := \langle \lambda, X_i(\pi(\lambda)) \rangle = 0,
\]
for some measurable controls \(u_i(t)\). An *abnormal extremal* is a nontrivial curve \(\lambda(t)\) in \(\Delta^\perp\) satisfying this system for some controls \(u(t)\).

Assume there exists a smooth, embedded submanifold \(A \subset \Delta^\perp\) that is invariant under all abnormal Hamiltonian flows (for all choices of controls) and such that the projection \(\pi|_A: A \to M\) is a submersion whose fibers \(A_x := \pi^{-1}(x)\) are compact and connected.

Define the set of projected tangent vectors
\[
E_x := \{\,d\pi(\xi) \mid \xi \in T_\lambda A, \; \xi \text{ tangent at } \lambda \in A_x \text{ to the abnormal Hamiltonian vector field}\,\}.
\]
Assume that \(\dim(E_x)\) is locally constant on \(M\), so \(E\) is a smooth distribution.

**Question:**  
1. Does the distribution \(E\) integrate to a (possibly singular) Riemannian foliation \(\mathcal{F}_E\) in the sense of Molino with respect to some Riemannian metric smoothly deformable from any extension of \(g\)?  
2. More specifically, does there exist a \(C^\infty\)-deformation \((\Delta_t, g_t)\), with \(\Delta_0 = \Delta\), such that abnormal geodesics of \((\Delta,g)\) are the \(t\to 0\) limits of normal geodesics of \((\Delta_t,g_t)\), and the limiting tangent distribution coincides with \(E\)?  

Why is it a "good" problem:  
This version clarifies the foundational setting of abnormal Hamiltonian dynamics precisely in Pontryagin’s framework, defines the objects through invariant geometric data, and asks whether the abnormal flow encodes a hidden integrable structure. It now rests on mathematically sound premises: a well-defined abnormal flow, a smooth invariant abnormal set, and a carefully constructed projected distribution.

The problem is **deep** because it touches several frontiers:  
1. **Control-theoretic singularities:** Understanding the smoothness and integrability of the abnormal set \(A\) would reveal when the endpoint mapping has structured singularities instead of chaotic ones.  
2. **Bridge to foliation theory:** Showing that \(E\) defines a Molino-type singular Riemannian foliation would transform our perception of abnormal extremals—from nonholonomic degeneracies to geometric leaves of a canonical foliation.  
3. **Deformation and rigidity:** The second question links the asymptotic behavior of geodesic flows under deformation to the stability and geometric “resolvability” of abnormals, which could lead to a notion of “regularization” for control singularities.  

The **simplicity** of its statement—asking whether invariant abnormal directions on \(T^*M\) induce geometric foliations on \(M\) and persist under deformation—belies the **depth** of analysis required, which spans symplectic geometry, geometric control, and foliation theory. It therefore fulfills the criteria of a *good* research-level problem: structurally clear, conceptually unbounded, and likely to open new mathematical connections rather than simply extend known ones.